%!TEX TS-program = xelatex
%!TEX TS-options = -shell-escape
%!TEX root = ../OpAlg2_SS16/operatoralgebren2.tex
\RequirePackage{fix-cm} 
\documentclass[a4paper, twoside, headsepline, index=totoc,toc=listof,toc=bibliography,toc=index, fontsize=10pt, cleardoublepage=empty, headinclude, DIV=12, BCOR=5mm, titlepage,draft]{scrartcl}
%!TEX root = ../AnaTopGeo_SS14/ana_top_geo.tex
\usepackage{scrtime} % KOMA, Uhrzeit ermoeglicht

%--Pakete zum "Programmieren"
% ======================================================================================
\usepackage{etoolbox}
\usepackage{letltxmacro}
\usepackage{ifthen}
% ======================================================================================

%--Farbdefinitionen und Grafiken (muss vor tikz geladen werden)
% ======================================================================================
\usepackage[usenames, table, x11names]{xcolor}
\definecolor{dark_gray}{gray}{0.45}
\definecolor{light_gray}{gray}{0.6}
\definecolor{fb10_blue}{cmyk}{0.8,0.4,0.13,0.07}
\usepackage[final]{graphicx}
\usepackage{adjustbox}
\newcommand{\cfbox}[2]{% coloured frame box
	\ifmmode
	\mathchoice{\adjustbox{cfbox=#1}{$\displaystyle#2$}}{\adjustbox{cfbox=#1}{$\textstyle#2$}}{\adjustbox{cfbox=#1}{$\scriptstyle#2$}}{\adjustbox{cfbox=#1}{$\scriptscriptstyle#2$}}
	\else
	\adjustbox{cfbox=#1}{#2}
	\fi
}
% ======================================================================================

%--Zum Zeichnen/ TikZ-Kram (vor polyglossia bzw. babel geladen werden)
% ======================================================================================
\usepackage{tikz}
\usepackage{tikz-cd}
\usetikzlibrary{external}
\tikzset{>=latex}
\usetikzlibrary{%
	shapes,
	arrows.meta,
	intersections,
	calc,
	3d,
	decorations.pathreplacing,decorations.markings,decorations.pathmorphing,
	angles,
	quotes,
}
\tikzexternalize[prefix=tikz/,up to date check=diff]
\pgfkeys{/pgf/images/include external/.code=\includegraphics{#1}}
\tikzset{external/system call={lualatex \tikzexternalcheckshellescape -halt-on-error -interaction=batchmode --shell-escape -jobname "\image" "\texsource"}}
\AtBeginEnvironment{tikzcd}{\tikzexternaldisable} % tikzexternalize fuer tikzcd deaktivieren, da inkompatibel
\AtEndEnvironment{tikzcd}{\tikzexternalenable}
\tikzset{% um Inkompatibilitaeten von quotes und polyglossia bzw. babel zu vermeiden
  every picture/.append style={
    execute at begin picture={\shorthandoff{"}},
    execute at end picture={\shorthandon{"}}
  }
}
\usepackage{pgfplots}
\usepgfplotslibrary{colormaps}
\newcommand*\circled[1]{\tikzexternaldisable\tikz[baseline=(char.base)]{\node[shape=circle,draw,inner sep=2pt] (char) {#1};}\tikzexternalenable}
% ======================================================================================



%-- Mathepakete etc.
% ======================================================================================
\usepackage[T1]{fontenc}
\renewcommand{\rmdefault}{zpltlf}
\usepackage{mathtools} % beinhaltet amsmath
\mathtoolsset{showonlyrefs,centercolon,showmanualtags}
\newtagform{brackets}[\textbf]{[}{]}
\usetagform{brackets}
\usepackage{fix-cm}
\usepackage[bbgreekl]{mathbbol}
\usepackage{amssymb,marvosym} 
\usepackage{nicefrac} % schräge Brüche
\usepackage{faktor}
\newcommand{\Faktor}[1]{\faktor[\textstyle]{#1}}
\usepackage{xfrac}
\usepackage{cancel}
\usepackage{mathdots} % Verbesserung von Punkten wie zB \ldots
\usepackage[bb=px]{mathalfa} % \mathbb als px font
\usepackage{centernot}
\usepackage{stackrel}
\DeclareSymbolFont{bbold}{U}{bbold}{m}{n}
\DeclareSymbolFontAlphabet{\mathbbold}{bbold}
\newcommand{\ind}{\mathbbold{1}} % charakteristische-Funktion-Eins
\def\mathul#1#2{\color{#1}\underline{{\color{black}#2}}\color{black}} %farbiges Untersteichen im Mathe-Modus
\renewcommand{\le}{\leqslant}
\renewcommand{\ge}{\geqslant}
% ======================================================================================


%-- Von xfrac erzeuge font warnings ignorieren
% ======================================================================================
\usepackage{silence}
\WarningFilter{latexfont}{Size substitutions with differences}
\WarningFilter{latexfont}{Font shape `U/bbold/m/n' in size}
% ======================================================================================


%-- Typographie/Polyglossia
% ======================================================================================
\usepackage[euler-digits]{eulervm} % vor fontspec laden!
\usepackage[no-math]{fontspec}
\usepackage{polyglossia} % moderner babel-ersatz
\setmainlanguage[spelling=new,babelshorthands=true]{german}
\shorthandoff{"}
\setotherlanguage{english}
\defaultfontfeatures{Mapping=tex-text, WordSpace={1.2}, Ligatures={Required,Common,Contextual},Extension=.otf} %


\setmainfont{TeXGyrePagellaX}[UprightFont=*-Regular,BoldFont=*-Bold,ItalicFont=*-Italic,BoldItalicFont=*-BoldItalic,ItalicFeatures={Style=Historic},Ligatures={Required,Common,Contextual,Historic}]
\setsansfont{texgyreadventor}[Scale=MatchUppercase, UprightFont=*-regular, BoldFont=*-bold, ItalicFont=*-italic, BoldItalicFont=*-bolditalic]
\setmonofont{SourceCodePro}[Scale=0.9,UprightFont=*-Regular, BoldFont=*-Semibold, ItalicFont=*-Light]
\usepackage{xltxtra}
\usepackage{fontawesome}
\usepackage[final]{microtype}
\usepackage[draft=false]{scrlayer-scrpage} 
\flushbottom
% ======================================================================================


%-- Aufzählungen
% ======================================================================================
\usepackage[shortlabels,inline]{enumitem}
\setlist[itemize,1]{label=\faCaretRight}
\setlist[enumerate]{font=\bfseries}
\setlist[description]{font=\normalfont\bfseries}
\usepackage{multicol}
% ======================================================================================


%-- Floats/Figures/Tabellen
% ======================================================================================
\usepackage{wrapfig}
\usepackage{float}
\usepackage[margin=10pt, font=small, labelfont={sf, bf}, format=plain, indention=1em]{caption}
\captionsetup[wrapfigure]{name=Abb. }
\usepackage{booktabs}
% ======================================================================================


%-- korrekte Anführungszeichen und Zitierbefehle
% ======================================================================================
\usepackage[autostyle,german=quotes,english=british]{csquotes}
% ======================================================================================


%--Indexverarbeitung
% ======================================================================================
\usepackage{makeidx}
\newcommand{\bet}[1]{\textbf{\emph{#1}}}
\newcommand{\Index}[1]{\bet{#1}\index{#1}}
\makeindex
\setindexpreamble{{\noindent\sffamily\small Die \emph{Seitenzahlen} sind mit Hyperlinks versehen und somit anklickbar} \par \bigskip}
\renewcommand{\indexpagestyle}{scrheadings}
% ======================================================================================


%-- Marginnotes/Todonotes/Footnotes
% ======================================================================================
\deffootnote[1.5em]{1.5em}{1.5em}{\textsuperscript{\thefootnotemark}\ }
\usepackage[fulladjust]{marginnote}
\renewcommand*{\marginfont}{\itshape\footnotesize}
\usepackage[textsize=small]{todonotes}
\usepackage{ragged2e}
\renewcommand*{\raggedleftmarginnote}{\RaggedLeft}
\renewcommand*{\raggedrightmarginnote}{\RaggedRight}
\LetLtxMacro{\oldtodo}{\todo}
\renewcommand{\todo}[2][]{\tikzexternaldisable\oldtodo[#1]{#2}\tikzexternalenable}
\LetLtxMacro{\oldmissingfigure}{\missingfigure}
\renewcommand{\missingfigure}[2][]{\tikzexternaldisable\oldmissingfigure[{#1}]{#2}\tikzexternalenable}
% ======================================================================================


% -- BibLaTeX
% ======================================================================================
\usepackage[%
	backend=biber,
	sortlocale=auto,
	natbib,
	hyperref,
	backref,
	style=alphabetic
	]%
{biblatex}
\renewcommand*{\mkbibnamelast}[1]{%
  \ifmknamesc{\textsc{#1}}{#1}}
\renewcommand*{\mkbibnameprefix}[1]{%
  \ifboolexpr{ test {\ifmknamesc} and test {\ifuseprefix} }
    {\textsc{#1}}
    {#1}}
\def\ifmknamesc{%
  \ifboolexpr{ test {\ifcurrentname{labelname}}
               or test {\ifcurrentname{author}}
               or ( test {\ifnameundef{author}} and test {\ifcurrentname{editor}} ) }}
\addbibresource{../!config/quellen.bib}
% ======================================================================================

%--Konfiguration von Hyperref und Cleveref
% ======================================================================================
\usepackage[hidelinks, pdfpagelabels,  bookmarksopen=true, bookmarksnumbered=true, linkcolor=black, urlcolor=SkyBlue2, plainpages=false,pagebackref, citecolor=black, hypertexnames=true, pdfauthor={Jannes Bantje}, pdfborderstyle={/S/U}, linkbordercolor=SkyBlue2, colorlinks=false,final,backref=false]{hyperref}
\usepackage[nameinlink,noabbrev]{cleveref}
\newcommand{\appendLink}[1]{#1\,\faExternalLink}
\newcommand{\hrefsym}[2]{\href{#1}{\texttt{\appendLink{#2}}}}
\newcommand{\hrefsymX}[2]{\href{#1}{\appendLink{#2}}}
\newcommand{\hrefsymmail}[2]{\href{#1}{\texttt{\faEnvelopeO\,#2}}}
\renewcommand{\url}[1]{\hrefsym{#1}{\nolinkurl{#1}}}
% ======================================================================================


% -- QR-Codes (hinter hyperref laden!)
% ======================================================================================
\usepackage{qrcode}
% ======================================================================================

%--Römische Zahlen
% ======================================================================================
\newcommand{\RM}[1]{\MakeUppercase{\romannumeral #1{}}}
% ======================================================================================

%-- Definition von diversen Mathe-Befehlen
% ======================================================================================
%!TEX root = mitschrift_main.tex

% -- Zum Finetuning von Befehlen
% ======================================================================================
\makeatletter
\newcommand{\raisemath}[1]{\mathpalette{\raisem@th{#1}}}
\newcommand{\raisem@th}[3]{\raisebox{#1}{$#2#3$}}
\makeatother
\makeatletter
\newcommand{\killDescendersM}[1]{\mathpalette{\killD@scendersM{#1}}}
\newcommand{\killD@scendersM}[2]{\raisebox{0pt}[\height][0pt]{$#2#1$}}
\makeatother
\DeclareRobustCommand{\minwidthbox}[2]{%
  \ifmmode
    \expandafter\mathmakebox
  \else
    \expandafter\makebox
  \fi
  [\ifdim#2<\width\width\else#2\fi]{#1}%
}
% ======================================================================================


%-- Klammerbefehle
% ======================================================================================
\DeclarePairedDelimiter{\abs}{\lvert}{\rvert}
\DeclarePairedDelimiter{\floor}{\lfloor}{\rfloor}
\DeclarePairedDelimiter{\ceil}{\lceil}{\rceil}
\DeclarePairedDelimiter\norm{\Vert}{\Vert}
\DeclarePairedDelimiter\enbrace{(}{)}
\DeclarePairedDelimiter\benbrace{[}{]}
\DeclarePairedDelimiter\bbenbrace{[\![}{]\!]}
\DeclarePairedDelimiter\lenbrace{<}{>}
\DeclarePairedDelimiter\angbrace{\langle}{\rangle}
\newcommand{\ssbrace}[1]{{\scriptscriptstyle\enbrace{#1}}}
\newcommand{\ssbbrace}[1]{{\scriptscriptstyle\benbrace{#1}}}
% ======================================================================================

%-- Mengen
% ======================================================================================
\newcommand\SetSymbol[1][]{\nonscript\:#1\vert\allowbreak\nonscript\:\mathopen{}}
\providecommand\given{} % to make it exist
\DeclarePairedDelimiterX\set[1]\{\}{\renewcommand\given{\SetSymbol[\delimsize]}#1}
% ======================================================================================

%-- Skalarprodukt (3 Varianten) 
% ======================================================================================
\DeclarePairedDelimiterX\sprod[2]{\langle}{\rangle}{#1\,\delimsize\vert\,#2}
\DeclarePairedDelimiterX\skal[2]{\langle}{\rangle}{#1\,,\,#2}
\makeatletter
\DeclareFontFamily{OMX}{MnSymbolE}{}
\DeclareSymbolFont{MnLargeSymbols}{OMX}{MnSymbolE}{m}{n}
\SetSymbolFont{MnLargeSymbols}{bold}{OMX}{MnSymbolE}{b}{n}
\DeclareFontShape{OMX}{MnSymbolE}{m}{n}{
    <-6>  MnSymbolE5
   <6-7>  MnSymbolE6
   <7-8>  MnSymbolE7
   <8-9>  MnSymbolE8
   <9-10> MnSymbolE9
  <10-12> MnSymbolE10
  <12->   MnSymbolE12
}{}
\DeclareFontShape{OMX}{MnSymbolE}{b}{n}{
    <-6>  MnSymbolE-Bold5
   <6-7>  MnSymbolE-Bold6
   <7-8>  MnSymbolE-Bold7
   <8-9>  MnSymbolE-Bold8
   <9-10> MnSymbolE-Bold9
  <10-12> MnSymbolE-Bold10
  <12->   MnSymbolE-Bold12
}{}
\let\llangle\@undefined
\let\rrangle\@undefined
\DeclareMathDelimiter{\llangle}{\mathopen}%
                     {MnLargeSymbols}{'164}{MnLargeSymbols}{'164}
\DeclareMathDelimiter{\rrangle}{\mathclose}%
                     {MnLargeSymbols}{'171}{MnLargeSymbols}{'171}
\makeatother
\DeclarePairedDelimiterX\sskal[2]{\llangle}{\rrangle}{#1\,,\,#2}
% ======================================================================================

%-- Abbildungsdefinition
% ======================================================================================
\newcommand{\mapdef}[5]{%
	\[
		\begin{array}{rcl}
			\textstyle #1 &\xrightarrow{\minwidthbox{#5}{2em}} & \textstyle #2 \\[0.5ex]
			\textstyle #3 &\xmapsto{\minwidthbox{\mbox{ }}{2em}} & \textstyle #4
		\end{array}
	\]
}
% ======================================================================================

%-- modifiziertes Stackrel 
% ======================================================================================
\newcommand{\StackText}[2]{\stackrel{\mbox{\scriptsize #1}}{#2}}
\newcommand{\StackTextClap}[2]{\stackrel{\mathclap{\mbox{\scriptsize #1}}}{#2}}
% ======================================================================================

%-- Blitz
% ======================================================================================
\newcommand{\light}{\text{\raisebox{-.3ex}{\Large\Lightning}}}
% ======================================================================================


%-- Underbrace u.Ä. als Befehl in LaTeX-Syntax (und ohne Spacingprobleme mit nachfolgenden Operatoren...)
% ======================================================================================
\newcommand{\Underbrace}[2]{{\underbrace{#1}_{#2}}}
\newcommand{\Underbracket}[2]{{\underbracket[0.7pt][2pt]{#1}_{#2}}}
\newcommand{\Overbracket}[2]{{\overbracket[0.7pt][2pt]{#1}^{#2}}}
% ======================================================================================


%-- Deklaration weiterer Operatoren (allgemein)
% ======================================================================================
\DeclareMathOperator{\re}{Re} % Realteil
\let\Re\relax
\DeclareMathOperator{\Re}{Re} % Realteil
\DeclareMathOperator{\im}{im} % Bild
\let\Im\relax
\DeclareMathOperator{\Im}{Im} % Bild
\DeclareMathOperator{\id}{id} % identische Abbildung
\DeclareMathOperator{\conj}{conj} % Konjugation
\DeclareMathOperator{\sgn}{sgn} % Signum
\DeclareMathOperator{\End}{End} % Endomorphismen
\DeclareMathOperator{\Hom}{Hom} % Homomorphismen
\DeclareMathOperator{\Iso}{Iso} % Isomorphismen
\DeclareMathOperator{\Aut}{Aut} % Automorphismen
\DeclareMathOperator{\Span}{span} % Span
\DeclareMathOperator{\coker}{coker} % Kokern
\DeclareMathOperator{\Tr}{Tr} % Spur,Trace
\DeclareMathOperator{\pr}{pr} % Projektion
\DeclareMathOperator{\diag}{diag} % Diagonalmatrix
\DeclareMathOperator{\Rg}{Rg} % Rang
\DeclareMathOperator{\const}{const} % konstante Abbildung
\DeclareMathOperator{\Spur}{Spur} % Spur
\DeclareMathOperator{\Arg}{Arg} % Argument
\DeclareMathOperator{\dist}{dist} % Distanz
\DeclareMathOperator{\supp}{supp} % Träger
\DeclareMathOperator{\Char}{char} % Charakteristik
% ======================================================================================


%-- Deklaration weiterer Operatoren (Differentiale etc.)
% ======================================================================================
\DeclareMathOperator{\grad}{grad} % Gradient
\DeclareMathOperator{\dive}{div} % Gradient
\DeclareMathOperator{\rot}{rot} % Rotation
\newcommand{\D}{\ensuremath{\mathrm{D}\mkern-1.0mu}} % Differential
\newcommand{\mathd}{\ensuremath{\mathrm{d}\mkern-1.0mu}} % äußere Ableitung
\newcommand{\Tmap}{\ensuremath{\mathrm{T}\mkern-0.85mu}} % Tangentialraum
\let\Tang\Tmap
\DeclareMathOperator{\Diff}{Diff}
\newcommand{\diff}[2]{\ensuremath{\frac{{\partial #1}}{{\partial #2}} }}
\newcommand{\diffd}[2]{\ensuremath{\frac{\mathd #1}{\mathd #2} }}
% ======================================================================================


%-- Deklaration weiterer Operatoren (Topologie)
% ======================================================================================
\newcommand*\interior[1]{\overset{\smash{\raisebox{-0.18ex}{$\scriptstyle\circ$}}}{#1}}
\newcommand{\sing}{{\raisemath{1.1pt}{\scriptscriptstyle\mathrm{sing}}}}
\newcommand{\pt}{\mathrm{pt}}
\DeclareMathOperator{\Zyl}{Zyl}
\DeclareMathOperator{\Tel}{Tel}
\newcommand{\op}{\mathrm{op}}
\DeclareMathOperator{\Sp}{Sp}
\DeclareMathOperator{\Keg}{Keg}
\newcommand{\slashedi}{i\hspace{-3.5pt}/}
\newcommand{\cupp}{\smallsmile}
\newcommand{\capp}{\smallfrown}
\DeclareMathOperator*{\colim}{colim}
\DeclareMathOperator{\PD}{PD}
\newcommand{\lf}{\mathrm{lf}}
\DeclareMathOperator{\sig}{sig}
\DeclareMathOperator{\Tor}{Tor}
\DeclareMathOperator{\Ext}{Ext}
\DeclareMathOperator{\AW}{AW}
\DeclareMathOperator{\Proj}{Proj}
\DeclareMathOperator{\Gr}{Gr}
\DeclareMathOperator{\res}{res}
\DeclareMathOperator{\Spec}{Spec}
\DeclareMathOperator{\co}{co}
\DeclareMathOperator{\ch}{ch}
\DeclareMathOperator{\wOp}{w}
\DeclareMathOperator{\Ar}{Ar}
\newcommand{\actson}{\mathrel{\curvearrowright}}
\let\acts\actson
\let\action\actson
\DeclareMathSymbol{\bbDelta}{\mathord}{bbold}{"01}
\newcommand{\DDelta}{\bbDelta}
\DeclareMathOperator{\Star}{Star}
\DeclareMathOperator{\Link}{Link}
\DeclareMathOperator{\EPK}{EPK}
\DeclareMathOperator{\Vol}{Vol}
\newcommand{\cell}{{\raisemath{1.1pt}{\scriptscriptstyle\mathrm{cell}}}}
\DeclarePairedDelimiter{\homologieklasse}{\llbracket}{\rrbracket}
\newcommand{\rand}[1]{\ensuremath{\partial^{\scriptscriptstyle #1}}}
\DeclareMathOperator{\ab}{ab}
% ======================================================================================


%-- Deklaration von Operatoren (Liegruppen)
% ======================================================================================
\DeclareMathOperator{\GL}{GL}
\DeclareMathOperator{\SO}{SO}
\DeclareMathOperator{\Ad}{Ad}
\DeclareMathOperator{\ad}{ad}
\DeclareMathOperator{\On}{O}
\DeclareMathOperator{\Un}{U}
\DeclareMathOperator{\SU}{SU}
\DeclareMathOperator{\Mat}{Mat}
\DeclareRobustCommand{\Der}{\mathop{\mathfrak{der}}}
\DeclareMathOperator{\SL}{SL}
\DeclareMathOperator{\Graph}{Graph}
\DeclareMathOperator{\Int}{Int}
\DeclareRobustCommand{\intAlg}{\mathop{\mathfrak{int}}}
\DeclareMathOperator{\aut}{aut}
\DeclareMathOperator{\Rad}{Rad}
\DeclareMathOperator{\Nil}{Nil}
\DeclareMathOperator{\rad}{rad}
\DeclareMathOperator{\nil}{nil}
\DeclareMathOperator{\Ric}{Ric}
\DeclareMathOperator{\ric}{ric}
\newcommand{\bi}{\mathrm{bi}}
\DeclareMathOperator{\Isom}{Isom}
\DeclareMathOperator{\Sym}{Sym}
\newcommand{\opL}{\ensuremath{\mathrm{L}\mkern-0.6mu}}
% ======================================================================================

%-- Deklaration von Operatoren (Funktionalanalysis)
% ======================================================================================
\DeclareMathOperator{\tr}{tr}
\newcommand{\w}{\mkern1mu\mathrm{w}}
\newcommand{\sa}{\mathrm{sa}}
\newcommand{\vb}{\mathrm{v\mkern-2.5mu.b\mkern-1.5mu.}} % vollständig beschränkt
\newcommand{\so}{\mathrm{\mkern.3mu s\mkern-1.4mu.\mkern-.6mu o\mkern-1.7mu.}} % \newcommand{\so}{\mathrm{s.o.}}
\newcommand{\solim}{\so\text{-}\mkern-0.8mu\lim}
\newcommand{\wo}{\mathrm{w\mkern-3mu.\mkern-.4mu o\mkern-1.7mu.}}
\newcommand{\Top}[1]{\mathcal{T}_{\mkern-2.3mu #1}}
\newcommand{\weakT}[1]{\ensuremath{\mathcal{T}_{#1}^{\mkern+1.0mu\text{\raisebox{0.4ex}{$\mathrm{w}$}}}}}
\newcommand{\weakTstar}[1]{\ensuremath{\mathcal{T}_{#1}^{\mkern+1.0mu\text{\raisebox{0.4ex}{$\mathrm{w}$}}^*}}}
\newcommand{\TWeakStar}{\Top{\w^*}}
\newcommand{\TWeakOp}{\Top{\wo}}
\newcommand{\Tso}{\Top{\so}}
\newcommand{\finSub}{\subset\mkern-0.7mu \subset}
\DeclareMathOperator{\Inv}{Inv}
\newcommand{\simm}{{\hspace{-1.6pt}\raisemath{0.5pt}{\sim}}}
\newcommand{\plus}{{\hspace{-1.6pt}+}}
\DeclareMathOperator{\ev}{ev}
\DeclareMathOperator{\Alg}{Alg}
\DeclareMathOperator{\her}{her}
\newcommand{\subher}{\subset_{\her}}
\newcommand{\grenzw}[1]{\xrightarrow{\minwidthbox{#1}{1.4em}}}
\newcommand{\grenzwl}[1]{\xleftarrow{\minwidthbox{#1}{1.4em}}}
\newcommand{\grenzwIn}[1]{\grenzw{\raisemath{-2pt}{#1}}}
\newcommand{\MyTo}[1]{\tikzexternaldisable\mathbin{\tikz[baseline] \draw[-to,line width=.4pt] (0ex,0.94ex) -- (#1,0.94ex);}\tikzexternalenable}
\newcommand{\dlim}{%
    \mathchoice
      {\lim\limits_{\MyTo{4.2ex}}}% \displaystyle
      {\lim\limits_{\MyTo{2.8ex}}}% \textstyle
      {\lim\limits_{\MyTo{2.3ex}}}% \scriptstyle
      {\lim\limits_{\MyTo{2.3ex}}}% \scriptscriptstyle
}
\newcommand{\Dlim}{\killDescendersM{\dlim}}
\DeclareMathOperator{\sep}{sep}
\DeclareMathOperator{\diam}{diam}
\DeclareMathOperator{\conv}{conv}
\DeclareMathOperator{\Prim}{Prim}
\DeclareMathOperator{\hull}{hull}
\DeclareMathOperator{\red}{red}
\DeclarePairedDelimiterX\bra[1]{\langle}{\rvert}{#1\,}
\DeclarePairedDelimiterX\ket[1]{\lvert}{\rangle}{\,#1}
\DeclarePairedDelimiterX\bracket[2]{\langle}{\rangle}{#1\,\delimsize\vert\,#2}
\newcommand{\tensormax}{\mathbin{\otimes_{\max}}}
\newcommand{\tensormin}{\mathbin{\otimes_{\min}}}
\DeclareMathOperator{\Ped}{Ped}
\newcommand{\alg}{\mathrm{alg}}
\DeclareMathOperator{\CPC}{CPC}
\DeclareMathOperator{\CP}{CP}
\DeclareMathOperator{\UPC}{UPC}
\newcommand{\DeltaOp}{\mathbin{\Delta}}
\newcommand{\kernedP}{\mathcal{P}\mkern-2mu}
\newcommand{\Pinfty}{\kernedP_{\infty}}
\DeclareMathOperator{\Groth}{Groth}
\DeclareMathOperator{\rk}{rk}
\newcommand{\MvN}{\mathrm{MvN}}
% ======================================================================================

%-- Kategorien
% ======================================================================================
\DeclareMathOperator{\Mor}{Mor}
\DeclareMathOperator{\Obj}{Obj}
\newcommand{\TOP}{\textsc{Top}}
\newcommand{\HTOP}{\textsc{HTop}}
\newcommand{\VR}{\textsc{VR}}
\newcommand{\MOD}{\textsc{Mod}}
\newcommand{\Mod}[1]{#1\text{-}\MOD}
\newcommand{\MONOIDE}{\textsc{Monoide}}
\newcommand{\SET}{\textsc{Set}}
\newcommand{\MAN}{\textsc{Man}}
\newcommand{\GRUPPEN}{\textsc{Gruppen}}
\newcommand{\ABELGRUPPEN}{\textsc{Abel.Gruppen}}
\newcommand{\ABEL}{\textsc{Abel}}
\newcommand{\KAT}{\textsc{Kat}}
\newcommand{\FUN}{\textsc{Fun}}
\newcommand{\SIMP}{\textsc{Simp}}
\newcommand{\VEKT}{\textsc{Vekt}}
\newcommand{\CH}{\textsc{Ch}}
\newcommand{\CSTARUN}{C^*\text{-}\textsc{Alg}^{\raisemath{-2.5pt}{1}}}
\newcommand{\CSTAR}{C^*\text{-}\textsc{Alg}}
\newcommand{\AB}{\textsc{Ab}}
% ======================================================================================
% ======================================================================================



% -- theorem packages
% ======================================================================================
\usepackage{amsthm}
\usepackage{thmtools,thm-restate}
\usepackage{mdframed}
\renewcommand{\listtheoremname}{Übersicht aller Aussagen}
\usepackage{bookmark}
\bookmarksetup{open,numbered}
\makeatletter
\newcommand*{\theorembookmark}{%
  \bookmark[
    dest=\@currentHref,
    rellevel=1,
    keeplevel,
  ]{%
    \thmt@thmname\space\csname the\thmt@envname\endcsname
    \ifx\thmt@shortoptarg\@empty
    \else
      \space(\thmt@shortoptarg)%
    \fi
  }%
}   
\makeatother
% ======================================================================================

% -- Definition der einzelnen Theorem-Umgebungen
% ======================================================================================
\declaretheoremstyle[%
	headfont=\sffamily\bfseries,
	notefont=\normalfont\sffamily\scshape,
	bodyfont=\normalfont,
	headformat=\NUMBER\ \NAME\NOTE,
	headpunct=.,
	postheadspace=1em,
	spaceabove=15pt,spacebelow=10pt,
	shaded={bgcolor=gray!20},
	postheadhook=\theorembookmark]%
{mainstyle}
\declaretheoremstyle[%
	headfont=\sffamily\bfseries,
	notefont=\normalfont\sffamily\scshape,
	bodyfont=\normalfont,
	headformat=\NUMBER\ \NAME\NOTE,
	headpunct=.,
	postheadspace=1em,
	spaceabove=15pt,spacebelow=10pt,
	shaded={bgcolor=fb10_blue!20},
	postheadhook=\theorembookmark]%
{mainstyle_blue}
\declaretheoremstyle[%
	headfont=\sffamily\bfseries,
	notefont=\normalfont\sffamily\scshape,
	bodyfont=\normalfont,
	headformat=\NUMBER\ \NAME\NOTE,
	headpunct=.,
	postheadspace=1em,
	spaceabove=15pt,spacebelow=10pt,
	postheadhook=\theorembookmark]%
{mainstyle_unshaded}
\declaretheoremstyle[%
	headfont=\sffamily\bfseries,
	notefont=\normalfont\sffamily\scshape,
	bodyfont=\normalfont,
	headformat=\NUMBER\NAME\NOTE,
	headpunct=.,
	postheadspace=1em,
	spaceabove=15pt,spacebelow=10pt,
	% shaded={bgcolor=gray!20},
	postheadhook=\theorembookmark]%
{mainstyle_unnumbered}
\declaretheoremstyle[%
	headfont=\sffamily\bfseries,
	notefont=\normalfont\sffamily\scshape,
	bodyfont=\normalfont,
	headformat=swapnumber,
	headpunct=.,
	postheadspace=1em,
	spaceabove=15pt,spacebelow=10pt,
	shaded={bgcolor=gray!20},
	postheadhook=\theorembookmark,
	qed=\qedsymbol]%
{mainstyleB}
\declaretheoremstyle[%
	headfont=\bfseries\scshape,
	bodyfont=\normalfont,
	headpunct=:,
	postheadspace=1em,
	spacebelow=12pt,spaceabove=2pt,
	qed=\qedsymbol]%
{beweise}
\declaretheoremstyle[%
	headfont=\bfseries\scshape,
	bodyfont=\normalfont,
	headpunct=:,
	postheadspace=1em,
	spacebelow=12pt,spaceabove=2pt]%
{beweisskizze}
\declaretheoremstyle[%
	headfont=\sffamily\bfseries,
	bodyfont=\normalfont,
	headpunct=:,
	postheadspace=1em,
	spacebelow=10pt,spaceabove=10pt]%
{bemerkungen}
\declaretheorem[name=Definition,parent=section,style=mainstyle_blue]{definition}
\declaretheorem[name=Definition \& Proposition,refname=Proposition,sharenumber=definition,style=mainstyle_blue]{definitionP}
\declaretheorem[name=Definition,numbered=no,style=mainstyle_unnumbered]{definition*}
\declaretheorem[name=Theorem,sharenumber=definition,style=mainstyle]{theorem}
\declaretheorem[name=Theorem,numbered=no,style=mainstyle_unnumbered]{theorem*}
\declaretheorem[name=Proposition,sharenumber=definition,style=mainstyle,refname=Proposition]{proposition}
\declaretheorem[name=Lemma,sharenumber=definition,style=mainstyle]{lemma}
\declaretheorem[name=Satz,sharenumber=definition,style=mainstyle,refname=Satz]{satz}
\declaretheorem[name=Satz,sharenumber=definition,style=mainstyle_unshaded]{satzUnshaded}
\declaretheorem[name=Definition,sharenumber=definition,style=mainstyle_unshaded]{definitionUnshaded}
\declaretheorem[name=Satz,numbered=no,style=mainstyle_unnumbered]{satz*}
\declaretheorem[name=Korollar,sharenumber=definition,style=mainstyle,refname=Korollar]{korollar}
\declaretheorem[name=Korollar,sharenumber=definition,style=mainstyleB,refname=Korollar]{korollarB}
\declaretheorem[name=Frage,numbered=no,style=mainstyle_unnumbered]{frage}
\declaretheorem[name=Frage,sharenumber=definition,style=mainstyle_unshaded]{frageA}
\declaretheorem[name=Erinnerung,sharenumber=definition,style=mainstyle_unshaded]{erinnerungA}
\declaretheorem[name=Ausblick,sharenumber=definition,style=mainstyle_unshaded]{ausblick}
\declaretheorem[name=Konvention,sharenumber=definition,style=mainstyle]{konvention}
\declaretheorem[name=Notation,sharenumber=definition,style=mainstyle_unshaded]{notation}
\declaretheorem[name=Bemerkung,sharenumber=definition,style=mainstyle_unshaded,refname=Bemerkung]{bemerkung}
\declaretheorem[name=Bemerkung,numbered=no,style=mainstyle_unnumbered]{bemerkung*}
\declaretheorem[name=Beispiel,sharenumber=definition,style=mainstyle_unshaded,refname=Beispiel]{beispiel}
\declaretheorem[name=Beispiel,numbered=no,style=mainstyle_unnumbered]{beispiel*}
\declaretheorem[name=Exkurs,numbered=no,style=mainstyle_unnumbered]{exkurs*}
\declaretheorem[name=Beweis,numbered=no,style=beweise]{beweis}
\declaretheorem[name=Übung,numbered=no,style=bemerkungen]{uebung}
\declaretheorem[name=Erinnerung,numbered=no,style=bemerkungen]{erinnerung}

% english versions
\declaretheorem[name=Remark,sharenumber=definition,style=mainstyle_unshaded]{remark}
\declaretheorem[name=Remark,numbered=no,style=mainstyle_unnumbered]{remark*}
\declaretheorem[name=Example,sharenumber=definition,style=mainstyle_unshaded]{example}
\declaretheorem[name=Corollary,sharenumber=definition,style=mainstyle]{corollary}
\let\proof\relax
\declaretheorem[name=Proof,numbered=no,style=beweise]{proof}
\declaretheorem[name=Sketch of Proof,numbered=no,style=beweisskizze]{sketch}
% ======================================================================================

%--Inhaltsverzeichnis
% ======================================================================================
\usepackage[tocindentauto]{tocstyle}
\usetocstyle{KOMAlike}
% ======================================================================================

%-- Dinge, die erst am Ende getan werden dürfen
% ======================================================================================
\shorthandon{"}
\usepackage{ellipsis}
% ======================================================================================


\newcommand{\fach}{Operatoralgebren \RM{2}.}
\newcommand{\semester}{Sose 2016}
\newcommand{\homepage}{http://wwwmath.uni-muenster.de/u/wilhelm.winter/wwinter/operatoralgebren_II.html}

\newcommand{\prof}{Prof.\ Dr.\ Wilhelm Winter}
\publishers{\scalebox{10}{\Huge$^{\raisebox{2.2pt}{\text{\small\RM{2}.}}}\mkern-1.5mu C^*$}}
\providecommand{\verfasser}{Jannes Bantje}

%--Konfiguration von scrheadings
\setheadsepline{1pt}[\color{light_gray}]
\pagestyle{scrheadings}
%\fancyhf[H, F]{}
\clearscrheadfoot

\providecommand{\shortFach}{\fach}
\lehead{\includegraphics[height=0.6 cm,keepaspectratio]{../!config/Bilder/Logo_WWU_Muenster_light_gray.pdf}}
\rehead{\rule{0cm}{0.6cm}\footnotesize\sffamily\color{light_gray}\verfasser{} -- Mitschrift \shortFach}
\lohead{\rule{0cm}{0.6cm}\footnotesize\sffamily\color{light_gray}Stand: \today \; \thistime[:]}
\rohead{\includegraphics[height=0.6 cm,keepaspectratio]{../!config/Bilder/fb10logo_gray.pdf}}


\ofoot[{ \color{dark_gray} \LARGE \sffamily \thepage}]{{ \color{dark_gray} \LARGE \sffamily \thepage}} %hier wir auch der plain Stil bearbeitet!
\automark{section}
\ifoot{ \color{dark_gray} \small \leftmark}

%--Metadaten
\providecommand{\mail}{j.bantje@wwu.de}
\author{\verfasser}
\titlehead{\includegraphics[height=1.5cm, keepaspectratio]{../!config/Bilder/Logo_WWU_Muenster.pdf}%
\hfill \includegraphics[height=1.3cm, keepaspectratio]{../!config/Bilder/fb10logo.pdf}}
\title{Skript \fach}
\subtitle{Mitschrift der Vorlesung  \enquote{\fach} von \prof}

\KOMAoptions{DIV=last}


\begin{document}
\pagenumbering{Roman}
\maketitle
\begin{abstract}
\section*{Aktuelle Version verfügbar bei}
\newcommand{\dieBreite}{11cm}
\begin{minipage}{4cm}
	\qrcode[height=3.3cm, version=6]{https://gitlab.com/JaMeZ-B/LaTeX-WWU}
\end{minipage}
\hfill
\begin{minipage}{\dieBreite}
	% \includegraphics[height=0.6cm, keepaspectratio]{../!config/Bilder/wm_no_bg.pdf}
	\includegraphics[height=0.8cm, keepaspectratio]{../!config/Bilder/wm_no_bg.pdf}\\
	\url{https://gitlab.com/JaMeZ-B/LaTeX-WWU} \smallskip\\
	Das zentrale Repository des \enquote{\LaTeX-WWU}-Projekts befindet sich auf der Plattform GitLab.com.
	Neben der Koordination aller Beteiligten werden über diesen Dienst auch die PDFs gebaut, die in der Readme verlinkt sind.
\end{minipage}\\[1cm]
\begin{minipage}{4cm}
	\qrcode[height=3.3cm, version=6]{https://github.com/JaMeZ-B/latex-wwu}
\end{minipage}
\hfill
\begin{minipage}{\dieBreite}
	\includegraphics[height=0.6cm, keepaspectratio]{../!config/Bilder/github_octo.pdf}
	\includegraphics[height=0.6cm, keepaspectratio]{../!config/Bilder/GitHub_Logo.pdf}\\
	\url{https://github.com/JaMeZ-B/latex-wwu} \smallskip\\
	Die Entwicklung des \enquote{\LaTeX-WWU}-Projekts hat ursprünglich auf GitHub stattgefunden, ist mittlerweile aber zu GitLab gewechselt.
	Das GitHub-Repository wird stündlich automatisch aktualisiert, Merge-Requests werden aber nicht mehr entgegengenommen.
\end{minipage}\\[1cm]
% \begin{minipage}{4cm}
% 	\qrcode[height=3.3cm, version=6]{https://uni-muenster.sciebo.de/public.php?service=files&t=965ae79080a473eb5b6d927d7d8b0462}
% \end{minipage}
% \hfill
% \begin{minipage}{\dieBreite}
% 	\raisebox{-2pt}{\includegraphics[height=0.6cm, keepaspectratio]{../!config/Bilder/sciebo_logo.pdf}}
% 	\resizebox{!}{0.5cm}{\large \sffamily\textbf{sciebo}} {\sffamily\large die Campuscloud} \\
% 	\resizebox{\dieBreite}{!}{\footnotesize\url{https://uni-muenster.sciebo.de/public.php?service=files&t=965ae79080a473eb5b6d927d7d8b0462}}\smallskip\\
% 	Sciebo ist ein Dropbox-Ersatz der Hochschulen in NRW, der von der Uni Münster in leitender Position auf Basis der OpenSource-Software Owncloud aufgebaut wurde.
% \end{minipage}\\[1cm]
\hrule \mbox{ }\\[0.7cm]
\begin{minipage}{4cm}
	\qrcode[height=3.3cm, version=6]{\homepage}
\end{minipage}
\hfill
\begin{minipage}{\dieBreite}
	\resizebox{!}{0.5cm}{\large\sffamily\textbf{Vorlesungshomepage}}\\
	\resizebox{\dieBreite}{!}{\footnotesize\url{\homepage}}\smallskip\\
	Hier ist ein Link zur offiziellen Vorlesungshomepage.
\end{minipage}
\newpage
\section*{Vorwort --- Mitarbeit am Skript}
Dieses Dokument ist eine Mitschrift aus der Vorlesung \enquote{\fach, \semester}, gelesen von \prof. 
Der Inhalt entspricht weitestgehend dem Tafelanschrieb. 
Für die Korrektheit des Inhalts übernehme ich keinerlei Garantie! 
Für Bemerkungen und Korrekturen -- und seien es nur Rechtschreibfehler -- bin ich sehr dankbar. 
Korrekturen lassen sich prinzipiell auf drei Wegen einreichen: 
\begin{itemize}
	\item Persönliches Ansprechen in der Uni, Mails an \hrefsymmail{mailto:\mail}{\mail} (gerne auch mit annotieren PDFs) oder Kommentare auf \url{https://gitlab.com/JaMeZ-B/LaTeX-WWU}.
	\item \emph{Direktes} Mitarbeiten am Skript: Den Quellcode poste ich auf GitLab (siehe oben), also stehen vielfältige Möglichkeiten der Zusammenarbeit zur Verfügung:
	Zum Beispiel durch Kommentare am Code über die Website und die Kombination Fork und Merge-Request. 
	Wer sich verdient macht oder ein Skript zu einer Vorlesung, die ich nicht besuche, beisteuern will, dem gewähre ich gerne auch Schreibzugriff.
	
	Beachten sollte man dabei, dass dazu ein Account bei \url{gitlab.com} notwendig ist, der allerdings ohne Angabe von persönlichen Daten angelegt werden kann. 
	Wer bei GitLab (bzw. dem zugrunde liegenden Open-Source-Programm \enquote{\texttt{git}}) -- verständlicherweise -- Hilfe beim Einstieg braucht, dem helfe ich gerne weiter. 
	Es gibt aber auch zahlreiche empfehlenswerte Tutorials im Internet.\footnote{zB. \url{https://try.github.io/levels/1/challenges/1}, ist auf Englisch, aber dafür interaktiv}
	\item \emph{Indirektes} Mitarbeiten: \TeX-Dateien per Mail verschicken. 
	
	Dies ist nur dann sinnvoll, wenn man einen ganzen Abschnitt ändern möchte (zB. einen alternativen Beweis geben), da ich die Änderungen dann per Hand einbauen muss! Ich freue mich aber auch über solche Beiträge!
\end{itemize}
\section*{Anmerkung}
Innerhalb dieser Mitschrift wird man öfter den Ausdruck \enquote{Warum?} finden. Dies sind vom Dozenten bewusst weggelassene Details, die zu verstärktem Mitdenken beim Lesen 
animieren sollen. Oftmals sind dies schon aus vorherigen Semestern bekannte Sachverhalte. Nur an wenigen Stellen habe ich die fehlenden Details hinzugefügt.

\section*{Literatur}
\begin{itemize}[itemsep=0pt]
	\item G. J. \textsc{Murphy} \emph{C*-algebras and Operator Theory} \cite{Murphy}
	\item N. P. \textsc{Brown}, N. \textsc{Ozawa} \emph{C*-algebras and Finite-Dimensional Approximations} \cite{brown2008}
	\item M. \textsc{Rørdam}, F. \textsc{Larsen}, N. J. \textsc{Laustsen} \emph{An Introduction to K-theory for C*-Algebras} \cite{rordam2000}
	\item N. E. \textsc{Wegge-Olsen} \emph{K-theory and C*-algebras} \cite{wegge}
	\item B. \textsc{Blackadar} \emph{K-theory for Operator Algebras} \cite{blackadar1998}
	\item B. \textsc{Blackadar} \emph{Operator Algebras} \cite{blackadar2006}
	\item M. \textsc{Takesaki} \emph{Theory of Operator Algebras I.} \cite{takesaki2002}
	\item R.V. \textsc{Kadison}, J. R. \textsc{Ringrose} \emph{Fundamentals of the Theory of Operator Algebras Volumes I. \& II.} \cite{kadison1997-1} \& \cite{kadison1997-2}
	\item G.K. \textsc{Pedersen} \emph{C*-algebras and their Automorphism Groups} \cite{PedC*}
\end{itemize}
\end{abstract}

\tableofcontents
\cleardoubleoddemptypage

\pagenumbering{arabic}
\setcounter{page}{1}
\setcounter{footnote}{0}

\section{Tensorprodukte} % (fold)
\label{sec:1}

\begin{erinnerungA}[{name=[Tensorprodukt]}]
	Seien $V,W$ $\mathbb{K}$-Vektorräume.
	Sei $X$ ein weiterer $\mathbb{K}$-Vektorraum und $\phi \colon V\times W \to X$ bilinear.
	Dann heißt $(X,\varphi)$ \Index{Tensorprodukt} von $V$ und $W$, falls folgende universelle Eigenschaft gilt:
	\[
		\begin{tikzcd}[sep=large]
			V \times W \drar["\text{bil.}"'] \rar["\varphi"] & X \dar[dashed,"\exists ! \psi \text{ lin.} "] \\
			& U
		\end{tikzcd}
	\]
\end{erinnerungA}

\begin{definitionP}[label=propdef:12,{name=[{Konstruktion des algebraischen Tensorproduktes}]}]
	$(X,\varphi)$ wie oben existiert und ist eindeutig bis auf linearen Isomorphismus.
	Wir schreiben $V \odot W$ für $X$ und $v \otimes w$ für $\varphi(v,w)$.\marginnote{$\odot$ algebraisches Tensorprodukt}
\end{definitionP}
\begin{beweis}
	Sei $M$ der freie Vektorraum mit Basis $V \times W$, anders ausgedrückt
	\[
		M:= \set[\big]{f \colon V\times W \to \mathbb{K} \given \supp f \text{ endlich}}
	\]
	Sei nun 
	\[
		N := \Span \set*{
			\begin{array}{l}
				(v_1 +v_2,w) - \enbrace[\big]{(v_1,w) + (v_2,w)} \\
				(v,w_1+w_2) - \enbrace[\big]{(v,w_1) + (v,w_2)} \\
				(v,\lambda \cdot w) - \lambda \cdot (v,w), (\lambda \cdot v,w) - \lambda \cdot (v,w) 
			\end{array}
		\given 
		\begin{array}{l}
			v,v_1,v_2 \in V,\\
			w,w_1,w_2 \in W,\\
			\lambda \in \mathbb{K}
		\end{array}}
	\]
	Wir setzen $X := \sfrac{M}{N}$ und $\varphi(v,w) := (v,w)+N$.
	Der Rest ist eine einfache Übung.
\end{beweis}

\begin{proposition}[{name=[Tensorprodukt von Algebren]}]
	Seien $A,B$ $\mathbb{K}$-Algebren.
	Dann ist $A \odot B$ wieder eine $\mathbb{K}$-Algebra mit 
	\[
		(a \otimes b) (a' \otimes b') = aa' \otimes bb'
	\]
	Sind $A$ und $B$ $^*$-Algebren, dann ist $A \odot B$ wieder eine $^*$-Algebra mit 
	\[
		(a \otimes b)^* = a^* \otimes b^*
	\]
	Sind $A$ und $B$ kommutativ/unital, so auch $A \odot B$.
\end{proposition}
\begin{beweis}
	Betrachte $M$ wie in \autoref{propdef:12} und überprüfe, dass Multiplikation auf $\sfrac{M}{N}$ wohldefiniert und assoziativ ist.
	
	Alternativ: Die Elemente $(a,b) \in A \times B$ liefern eine bilineare Abbildung
	\mapdef{M_{(a,b)} \colon A \times B}{A \odot B}{(a',b')}{aa' \otimes bb'}{}
	Aus der universellen Eigenschaft erhält man $m_{(a,b)} \colon A \odot B \to A \odot B$ linear.
	Dies liefert wiederum eine bilineare Abbildung $m \colon A \times B \to \mathcal{L}(A \odot B, A \odot B)$ durch $(a,b) \mapsto m_{(a,b)}$.
	Nach der universellen Eigenschaft bekommt man nun eine lineare Abbildung $A \odot B \to \mathrm{L}(A \odot B, A\odot B)$.
	Daraus erhält man eine bilineare Abbildung
	\[
		\mu \colon A \odot B \times A \odot B \longrightarrow A \odot B
	\]
	Dann ist $(A \odot B,\mu)$ eine $\mathbb{K}$-Algebra.
\end{beweis}

\begin{bemerkung}[label=bem:14,{name=[{universelle Eigenschaft des Tensorproduktes von *-Algebren}]}]
	\begin{enumerate}[(i)]
		\item Seien $A,B$ $^*$-Algebren.
		Dann hat das Tensorprodukt $A \odot B$ als $^*$-Algebra die folgende universelle Eigenschaft:
		
		Falls $C$ eine $^*$-Algebra ist und $\pi_A \colon A \to C$, $\pi_B \colon B \to C$ $^*$-Homomorphismen sind, sodass für den Kommutator $\benbrace*{\pi_A(A), \pi_B(B)}=0$ gilt, so existiert genau ein $^*$-Homomorphismus $\pi \colon A \odot B \to C$ mit $\pi(a \otimes b)= \pi_A(a) \cdot \pi_B(b)$ für alle $a \in A$, $b \in B$.
		\item Es seien $A,B,C,D$ $^*$-Algebren, $\pi_A \colon A \to C$ und $\pi_B \colon B \to D$ $^*$-Homomorphismen.
		Dann existiert genau ein $\pi \colon A \odot B \to C \odot D$ mit $\pi(a \otimes b)= \pi_A(a) \otimes \pi_B(b)$.
	\end{enumerate}
\end{bemerkung}

\begin{definitionP}[{name=[{Tensorprodukt von Hilberträumen}]}]
	Es seien $\mathcal{H}_1$, $\mathcal{H}_2$ Hilberträume.
	Dann ist $\mathcal{H}_1 \odot \mathcal{H}_2$ ein Prä-Hilbertraum mit
	\[
		\skal*{\xi \otimes \eta}{\xi' \otimes \eta'} := \skal*{\xi}{\xi'} \cdot \skal*{\eta}{\eta'} 
	\]
	Das \bet{Tensorprodukt von Hilberträumen}\index{Tensorprodukt!von Hilberträumen} ist nun definiert durch
	\[
		\mathcal{H}_1 \otimes \mathcal{H}_2 := \overline{\mathcal{H}_1 \odot \mathcal{H}_2}^{\skal*{\cdot}{\cdot}}
	\]
\end{definitionP}
\begin{beweis}
	Es seien $\xi \in \mathcal{H}_1$ und $\eta \in \mathcal{H}_2$.
	Nach der universellen Eigenschaft existiert eine lineare Abbildung $\tau_{(\xi,\eta)} \colon \mathcal{H}_1 \odot \mathcal{H}_2 \to \mathbb{C}$ mit $\tau_{{\xi,\eta}} \enbrace*{\xi' \otimes \eta'} = \skal*{\xi}{\xi'} \cdot \skal*{\eta}{\eta'}$.
	Die Abbildung $(\xi,\eta) \mapsto \overline{\tau_{{\xi,\eta}}}$ ist dann bilinear.
	Aus der universellen Eigenschaft erhalten wir eine antilineare Abbildung $\xi \otimes \eta \mapsto \tau_{(\xi,\eta)}$.
	Daraus erhält man eine Sesquilinearform 
	\[
		\skal*{\cdot}{\cdot } \colon \mathcal{H}_1 \odot \mathcal{H}_2 \longrightarrow \mathbb{C}
	\]
	Wir müssen noch zeigen, dass diese positiv definit ist:
	Sei $z = \sum_{i=1}^{n} \xi_i \otimes \eta_i \in \mathcal{H}_1 \odot \mathcal{H}_2$.
	Wähle eine Orthonormalbasis $\zeta_1, \ldots , \zeta_m$ für $\Span \set*{\eta_1, \ldots,\eta_n}$.
	Dann ist $z=\sum_{j=1}^{m} \xi_j' \otimes \zeta_j$ für $\xi_j' \in \mathcal{H}_1$ geeignet.
	Damit folgt 
	\[
		\skal*{z}{z} = \sum_{j=1}^{m} \norm*{\xi'_j}^2 \ge 0
	\]
	Insbesondere ist $z=0$ genau dann, wenn $\xi_1' = \ldots = \xi_m'=0$ ist, also $\skal*{z}{z}=0$ gilt.
\end{beweis}

\begin{bemerkung}[{name=[{Einbettung der Tensorprodukte beschränkter Operatoren}]},label=bem:16]
	Es existiert eine Einbettung $\iota \colon \mathcal{B}(\mathcal{H}_1) \odot \mathcal{B}(\mathcal{H}_2) \hookrightarrow \mathcal{B}(\mathcal{H}_1 \otimes \mathcal{H}_2)$ gegeben durch $(a \otimes b)(\xi \otimes \eta) = a(\xi) \otimes b (\eta)$. (Warum?)
\end{bemerkung}

\begin{frageA}[{name=[{C*-Normen auf dem algebraischen Tensorprodukt}]}]
	Seien $A$ und $B$ $C^*$-Algebren.
	Die spannende Frage ist nun:
	\begin{center}
		\fbox{{\large¿}\,Was für $C^*$-Normen existieren auf $A \odot B$\,{\large?}}
	\end{center}
\end{frageA}

\begin{bemerkung}[{name=[Tensorprodukt aus C*-Norm auf dem algebraischen Tensorprodukt]}]
	Falls $\gamma \colon A \odot B \to \mathbb{R}^+$ eine $C^*$-Norm ist, so ist $A \otimes_\gamma B := \overline{A \odot B}^\gamma$ eine $C^*$-Algebra.\todo{RevChap 1}
\end{bemerkung}

\begin{proposition}[{name=[{minimale C*-Norm auf dem Tensorprodukt}]}]
	Es seien $A$ und $B$ $C^*$-Algebren.
	\begin{enumerate}[(i)]
		\item Für $^*$-Darstellungen $\pi_A \colon A \to \mathcal{B}(\mathcal{H}_1)$, $\pi_B \colon B \to \mathcal{B}(\mathcal{H}_2)$ definiert 
		\[
			(\pi_A \otimes \pi_B)(a \otimes b) := \iota\enbrace[\big]{\pi_A(a) \otimes \pi_B(b)}
		\]
		eine $^*$-Darstellung $\pi_A \otimes \pi_B \colon A \odot B \to \mathcal{B}(\mathcal{H}_1 \otimes \mathcal{H}_2)$.
		Falls $\pi_A$ und $\pi_B$ treu sind, so ist auch $\pi_A \otimes \pi_B$ treu.
		\item Die Abbildung $\norm*{\cdot }_{\min} \colon A \odot B \to \mathbb{R}^+$, gegeben durch
		\[
			\norm[\big]{\textstyle\sum\nolimits_{i=1}^{n} a_i \otimes b_i}_{\min} := \sup \set[\Big]{\norm[\big]{(\pi_A \otimes \pi_B)\enbrace[\big]{\textstyle\sum\nolimits_{i=1}^{n} a_i \otimes b_i}} \given \pi_A, \pi_B \text{ $^*$-Darstellungen} }
		\]
		ist eine $C^*$-Norm.
	\end{enumerate}
\end{proposition}
\begin{beweis}
	\begin{enumerate}[(i)]
		\item Nach \autoref{bem:14} und \autoref{bem:16} existiert genau eine $^*$-Darstellung 
		\[
			\pi_A \otimes \pi_B = \iota \circ (\pi_A \otimes \pi_B)
		\]
		Seien $\pi_A$ und $\pi_B$ treu und $0 \neq c \in \ker (\pi_A \otimes \pi_B)$.
		Es gilt $c= \sum_{i=1}^{n} a_i \otimes b_i$ für geeignete $a_i$, $b_i$ und $n$.
		Ohne Einschränkungen können wir annehmen, dass die $b_i$ linear unabhängig sind.
		Da $\pi_B$ treu ist, sind auch die $\pi_B(b_i)$ linear unabhängig.
		Für $\xi \in \mathcal{H}_1$ wählen eine Orthonormalbasis $\set*{\xi_1, \ldots ,\xi_m}$ für $\Span \set*{\pi_A(a_i) \xi \given i=1,\ldots ,n} \subset \mathcal{H}_1$.
		Es existieren also $\lambda_{ij} \in \mathbb{C}$ für $i=1,\ldots ,n$, $j=1,\ldots ,m$ mit 
		\(
			\pi_A(a_i) \xi = \sum_{j=1}^{m} \lambda_{ij} \cdot \xi_j
		\)
		für $i=1,\ldots ,n$.
		Für $\eta \in \mathcal{H}_2$ gilt
		\begin{align}
			0 = \enbrace*{\pi_A \otimes \pi_B}(c) \enbrace*{\xi \otimes \eta} = \sum_{i=1}^{n} \pi_A(a_i) \xi \otimes \pi_B(b_i) \eta &= \sum_{i=1}^{n} \enbrace*{\sum_{j=1}^{m} \lambda_{ij} \cdot \xi_j} \otimes \pi_B(b_i) \eta \\
			&= \sum_{j=1}^{m} \xi_j \otimes \enbrace*{\sum_{i=1}^{n} \lambda_{ij} \cdot \pi_B(b_i) \eta }
		\end{align}
		Es folgt, dass $\sum_{i=1}^{n} \lambda_{ij} \cdot \pi_B(b_i) \eta=0$ für $j=1,\ldots ,m$ und $\eta \in \mathcal{H}_2$.
		Damit sind die $\pi_B(b_i)$ linear abhängig oder $\lambda_{ij}=0$.
		Ersteres kann nicht gelten, also muss $\pi_A(a_i)\xi=0$ für $\xi \in \mathcal{H}_1$ gelten, also $\pi_A(a_i)=0$.
		Folglich ist $c=0$ im Widerspruch zur Annahme.
		\item $C^*$-Halbnorm: Übung!
		Wir zeigen, dass $\norm*{\cdot }_{\min}$ endlich ist: 
		\begin{align}
			\norm*{\enbrace*{\pi_A \otimes \pi_B} \enbrace*{\sum\nolimits_{i=1}^{n} a_i \otimes b_i}} &= \norm*{\sum_{i=1}^{n} \enbrace*{\pi_A(a_i) \otimes \ind_{\mathcal{H}_1} } \enbrace*{\ind_{\mathcal{H}_2} \otimes \pi_B(b_i)}}\\
			&\le \sum_{i=1}^{n} \norm*{\pi_A(a_i) \otimes \ind_{\mathcal{H}_1}} \norm*{\ind_{\mathcal{H}_2} \otimes \pi_B(b_i)} \\
			\intertext{$\pi_A(\cdot ) \otimes \ind_{\mathcal{H}_1}$ und $\ind_{\mathcal{H}_2} \otimes \pi_B(\cdot )$ sind als $^*$-Darstellungen kontraktiv, also}
			&\le \sum_{i=1}^{n} \norm*{a_i} \norm*{b_i} < \infty
		\end{align}
		$\norm*{\cdot }_{\min}$ ist eine Norm nach (i), denn $A$ und $B$ besitzen treue Darstellungen.\qedhere
	\end{enumerate}
\end{beweis}

\begin{definition}[{name=[{minimales Tensorprodukt}]}]
	\[
		\overline{A \odot B}^{\norm*{\cdot }_{\min}} =: A \otimes_{\min}B \enspace(= {A \otimes B})
	\]
	heißt \bet{minimales} (oder \bet{räumliches} oder \bet{spatiales}) \bet{Tensorprodukt}\index{minimales Tensorprodukt} von $A$ und $B$.
\end{definition}

\begin{satz}
	Seien $A,B$ $C^*$-Algebren und $\pi \colon A \odot B \to \mathcal{B}(\mathcal{H})$ eine nichtdegenerierte Darstellung.
	Dann existieren eindeutig bestimmte nichtdegenerierte Darstellungen $\pi_A \colon A \to \mathcal{B}(\mathcal{H})$ und $\pi_B \colon B \to \mathcal{B}(\mathcal{H})$ mit 
	\[
		\pi(a \otimes b) = \pi_A(a) \cdot \pi_B(b)=\pi_B(b) \cdot \pi_A(a)
	\]
	Falls $\pi$ eine Faktordarstellung ist (das heißt $\overline{\pi(A \odot B)}^{\wo}$ ist ein Faktor), so auch $\pi_A$ und $\pi_B$.
\end{satz}
\begin{beweis}
	Falls $A$ und $B$ unital sind, so sind $\pi_A$, $\pi_B$ gegeben durch $\pi_A(a) = \pi(a \otimes \ind_B)$ bzw. $\pi_B(b) = \pi(\ind_A \otimes b)$.
	Im allgemeinen Fall seien $(h_\lambda), (k_\mu)$ approximative Einsen.
	Setze 
	\begin{align}
		\pi_A(a) &:= \solim \pi(a \otimes k_\mu) \\
		\pi_B(b) &:= \solim \pi(h_\lambda \otimes b)
	\end{align}
	Wir zeigen, dass diese Limiten existieren:
	Für $\xi \in \mathcal{H}$ und $a \in A_+$ definiert $b \mapsto \skal*{\xi}{\pi(a \otimes b) \xi}$ ein positives Funktional $f_a$ auf $B$ (beschränkt nach \cite[Satz 5.3]{OpAlg}).
	Damit existiert
	\[
		\lim_\mu f_a(k_\mu) = \lim_\mu \skal*{\xi}{\pi(a \otimes k_\mu)\xi}
	\]
	Für $a \in A_+$, $\mu \le \mu'$ gilt 
	\begin{align}
		\norm[\Big]{\enbrace[\big]{\pi(a \otimes k_{\mu'}) - \pi (a \otimes k_\mu)}\xi}^2 = \skal[\Big]{\xi}{\pi \enbrace[\big]{a^2 \otimes \enbrace*{k_{\mu'} - k_\mu}^2} \xi}
		&\le \skal[\Big]{\xi}{ \pi \enbrace[\big]{a \otimes \enbrace*{k_{\mu'}- k_\mu}} \xi} \\
		&= f_a(k_{\mu'}) - f_a(k_\mu) < \varepsilon
	\end{align}
	für $\mu, \mu'$ groß genug.
	Also existiert $\pi_A(a)$ für $a \in A_+$. 
	Mit Linearität folgt, dass $\pi_A$ auch insgesamt wohldefiniert ist.
	Dass $\pi_A, \pi_B$ $^*$-Darstellungen sind mit den gewünschten Eigenschaften, verbleibt als Übung!
	
	Es gilt nach dem Bikommutantensatz
	\[
		\begin{tikzcd}
			\pi_A(A)\dar[phantom,"\subset", sloped] \rar[phantom,"\subset", sloped] & \pi_B(B)' \dar[equals] \\
			\pi_A(A)'' \rar[phantom,"\subset", sloped] & \pi_B(B)''
		\end{tikzcd}
	\]
	Damit folgt $\pi_A(A)' \cap \pi_A(A)'' \subset \pi_A(A)' \cap \pi_B(B)' \subset \pi(A \odot B)'$.
	Weiter haben wir $\pi_A(A) \subset \overline{\pi(A \odot B)}^{\so} = \pi(A \odot B)''$, also gilt $\pi_A(A)'' \subset \pi(A \odot B)''$, womit weiter folgt
	\[
		\mathbb{C} \cdot \ind_{\mathcal{H}} \subset \pi_A(A)' \cap \pi_A(A)'' \subset \pi(A \odot B)' \cap \pi(A \odot B)'' = \mathbb{C} \cdot \ind_{\mathcal{H}}\qedhere
	\]
\end{beweis}

\begin{korollar}
	Seien $A,B$ $C^*$-Algebren, $\gamma$ eine $C^*$-Halbnorm auf $A \odot B$.
	Dann gilt $\gamma(a \otimes b) \le \norm*{a} \cdot \norm*{b}$ für $a \in A$, $b	\in B$.
	$\gamma$ besitzt eine eindeutige Fortsetzung zu einer $C^*$-Halbnorm auf $A^\simm \odot B^\simm$.
\end{korollar}
\begin{beweis}
	Setze $N := \set*{x \in A \odot B \given \gamma(x)=0} \lhd A \odot B$.
	$\gamma$ induziert eine $C^*$-Norm $\overline{\gamma}$ auf $\sfrac{A \odot B}{N}$.
	Damit ist $C := \overline{\sfrac{A \odot B}{N}}^{\overline{\gamma}}$ eine $C^*$-Algebra.
	$C$ besitzt eine nichtdegenerierte treue Darstellung.
	Wir erhalten $\pi \colon A \odot B \to \mathcal{B}(\mathcal{H})$ mit $\norm*{\pi(X)}= \overline{\gamma}(x +N) =\gamma(x)$ für $x \in A \odot B$.
	Weiter gilt
	\[
		\gamma(a \otimes b) = \norm*{\pi(a \otimes b)} \le \norm*{\pi_A(a)} \cdot \norm*{\pi_B(b)} \le \norm*{a} \cdot \norm*{b}
	\]
	Seien $\tilde{\pi}_A \colon A^\simm \to \mathcal{B}(\mathcal{H})$, $\tilde{\pi}_B \colon B^\simm \to \mathcal{B}(\mathcal{H})$ die Fortsetzungen auf den Unitalisierungen.
	Dann definiert $\tilde{\pi}(a \otimes b) := \tilde{\pi}_A(a) \cdot \tilde{\pi}_B(b)$ eine Fortsetzung $\tilde{\pi} \colon A^\simm \odot B^\simm \to \mathcal{B}(\mathcal{H})$ von $\pi$.
	($\tilde{\pi}$ existiert nach \autoref{bem:14}).
	$\tilde{\gamma}(\cdot ) := \norm*{\tilde{\pi}(\cdot)}$ ist eine Fortsetzung von $\gamma$.
	Wir müssen die Eindeutigkeit zeigen:
	\begin{align}
		\tilde{\gamma}(x) = \norm*{\tilde{\pi}(x)} = \norm*{\ind_{\mathcal{H}} \cdot \tilde{\pi}(x)} &= \norm*{\solim \pi\enbrace[\big]{h_\lambda^{\sfrac{1}{2}} \otimes k_\mu^{\sfrac{1}{2}}} \cdot \tilde{\pi}(x)} \\
		&= \lim_{\lambda,\mu} \norm*{\tilde{\pi}(x^*) \pi \enbrace*{h_\lambda \otimes k_\mu} \tilde{\pi}(x)}^{\sfrac{1}{2}} \\
		&= \lim_{\lambda,\mu} \norm*{\pi \enbrace*{x^* (h_\lambda \otimes k_\mu) x}}^{\sfrac{1}{2}} \\
		&= \lim_{\lambda,\mu} \gamma \enbrace*{x^*(h_\lambda \otimes k_\mu)x}
	\end{align}
	Der letzte Ausdruck ist unabhängig von $\pi$.
\end{beweis}

% section 1 (end)

\cleardoubleoddemptypage
\pagenumbering{Alph}
\setcounter{page}{1}
\cleardoubleoddemptypage
\appendix

\section{Anhang} % (fold)
\label{sec:anhang}
%!TEX root = ana_top_geo.tex

\subsection{Ausführlicher Beweis zu \cref{lem:kpt-schnitte}} % (fold)
\label{sub:kpt-schnitte}
Sei $X$ ein Hausdorffraum. Dann ist $X$ genau dann kompakt, wenn gilt: Hat eine Familie $\mathcal{A}$ von abgeschlossenen Teilmengen von $X$ die endliche 
Durchschnittseigenschaft, so gilt 
\[
	\bigcap_{A \in \mathcal{A}} A \not= \emptyset.
\]
\begin{beweis}
	Für die erste Implikation sei $X$ kompakt und $\mathcal{A}$ eine Familie von abgeschlossenen Mengen mit der endlichen Durchschnittseigenschaft.
	Angenommen $\bigcap_{A \in \mathcal{A}} A = \emptyset$.
	Dann gilt
	\[
		X = X \setminus \bigcap_{A \in \mathcal{A}} A = \bigcup_{A \in \mathcal{A}} X \setminus A.
	\]
	Nun ist $\mathcal{U} \coloneqq \set*{X \setminus A \given A \in \mathcal{A}}$ eine offene Überdeckung von $X$ und da $X$ kompakt ist, existiert $\mathcal{A}_0 \subset \mathcal{A}$ endlich, sodass
	\[
		X = \bigcup_{A \in \mathcal{A}_0} X \setminus A = X \setminus \underbrace{\bigcap_{A \in \mathcal{A}_0 } A }_{\neq \emptyset} \quad \light
	\]
	Für die umgekehrte Implikation sei nun $\mathcal{U} = \set{U_i}_{i \in I}$ eine offene Überdeckung von $X$.
	Angenommen für jede endliche Teilmenge $J \subseteq I$ gilt $X \neq \bigcup_{i \in J} U_i$.
	Betrachte nun $\mathcal{A} =  \set{X \setminus U_i}_{i \in I}$. Dann gilt nach Annahme
	\[
		\bigcap_{i \in J} X \setminus U_i = X \setminus \bigcup_{i \in J} U_i \neq \emptyset.
	\]
	Also hat $\mathcal{A}$ die endliche Durchschnittseigenschaft. Nach Vorraussetzung gilt dann
	\[
		\emptyset \not= \bigcap_{i \in I} X \setminus U_i = X \setminus \underbrace{\bigcup_{i \in I} U_i}_{= X} \quad \light \qedhere
	\]
\end{beweis}


\printindex
\printbibliography
\listoffigures
\todototoc
\listoftodos[To-do's und andere Baustellen]
\end{document}
